% TOP_PROBLEM: {"id": 3170, "title": "Jones' conjecture", "category_id": 3, "category": {"id": 3, "name": "graph_theory", "display_name": "Graph Theory", "description": "Problems involving graphs, networks, and their properties.", "slug": "graph-theory", "order_index": 3, "created_at": "2026-07-31T15:26:25.671Z"}, "status": "open", "problem_number": "OPG-638", "set_id": 12}
% FIRST_POSTED: 2026-10-01
% ENTRY_KIND: statement_only
% UnsolvedMath ID 3170; source catalogue code OPG-638
% !TeX program = lualatex
% Definition and sources only; this file is not an LLM solution attempt.
\documentclass[11pt,a4paper]{article}
\usepackage[margin=25mm]{geometry}
\usepackage{fontspec}
\setmainfont{lmroman10-regular.otf}[BoldFont=lmroman10-bold.otf,
 ItalicFont=lmroman10-italic.otf,BoldItalicFont=lmroman10-bolditalic.otf]
\usepackage{amsmath,amssymb,mathtools}
\usepackage{unicode-math}
\setmathfont{latinmodern-math.otf}
\AtBeginDocument{\let\setminus\smallsetminus}
\usepackage{xurl,hyperref}
\hypersetup{colorlinks=true,urlcolor=blue}
\setcounter{secnumdepth}{0}
\setlength{\parindent}{0pt}
\setlength{\parskip}{5pt}
\setlength{\emergencystretch}{3em}
\begin{document}
\section*{Jones' conjecture}\label{3170}
{\small UnsolvedMath ID 3170 · OPG-638 · Graph Theory · OpenGarden}
\subsection{Definitions and mathematical statement}
Let $G=(V,E)$ be a finite simple planar graph, so it admits
an embedding in the plane without crossing edges. A cycle
means a connected $2$-regular subgraph. A cycle packing is
a collection of cycles with pairwise disjoint vertex sets;
write $cp(G)$ for the largest possible number of cycles in
such a collection. The quantity counts cycles, not their
total number of vertices or edges.

A feedback vertex set is a set $X\subseteq V$ such that
$G-X$ is acyclic, where deleting vertices also deletes all
their incident edges. Write $cc(G)$ for the smallest size
of a feedback vertex set. Equivalently, $X$ must meet the
vertex set of every cycle in $G$.

Jones' conjecture is the uniform inequality
\[
 cc(G)\leq2\,cp(G)
 \qquad\text{for every finite simple planar graph }G.
\]
Thus, if the largest family of vertex-disjoint cycles has
$k$ members, at most $2k$ vertices should suffice to hit all
cycles, including overlapping cycles outside a chosen packing.
The vertices of the hitting set may be anywhere in $G$;
they need not be selected only from one optimal packing.

For a forest both parameters are zero. Disconnected graphs
are allowed, with both parameters adding over components.
The elementary reverse bound $cp(G)\leq cc(G)$ follows
because disjoint cycles require distinct deleted vertices.
The proposed factor two is the additional conclusion
supplied by planarity.
\subsection{Short English statement}
In every planar graph, can all cycles be destroyed by deleting at most twice as many vertices as the maximum number of vertex-disjoint cycles?
\subsection{Sources}
\begin{itemize}
\item[S1] ULAM.AI and the UnsolvedMath Contributors, supplied problems.json, record 3170 (OPG-638), OpenGarden collection. Catalogue identity and source transcription; CC BY 4.0.
\url{https://huggingface.co/datasets/ulamai/UnsolvedMath}.
\item[S2] Open Problem Garden, Jones' conjecture. Original problem and discussion; the mathematical formulation below is expanded from the supplied source record.
\url{http://www.openproblemgarden.org/op/jones_conjecture}.
\end{itemize}
\end{document}
